\section{Shared activity and material-flow framework}\label{sec:architecture}
\subsection{Scope and canonical source records}
Each source record distinguishes physical source, emitting process, reporting category, calculation module, native species and chemical mass convention. Activity quantity, unit, fiscal dates and provenance accompany factor applicability, technology, control assumptions, configuration, estimate maturity, coverage and temporal/spatial support. Calculation-module totals are implementation summaries; category accounts determine scientific reporting and prevent captive fuel or process carbon from being counted twice.

For an appropriate activity/factor process,
\begin{equation}E_{s,p,t}=A_{s,t}EF_{s,p,t}(1-\eta_{s,p,t}).\end{equation}
This expression is not applied indiscriminately. Sulfur balances, surviving nitrogen pools, material conversion and recovered process carbon have their own equations. A factor that already describes controlled operation is not discounted again. Activity is shared where the physical stream is shared; pollutant-specific chemistry and process boundaries remain distinct.

\subsection{Fuel and industrial activity}
OCAC product-specific sales are reconciled at compatible use boundaries before road, off-road, industrial, railway and household fuel are partitioned. Overlapping editions identify historical revisions; a revised national total does not invent unpublished revised sector cells. OGRA deliveries, physical production and infrastructure distinguish gas use from fugitive pathways. Unaccounted-for gas is not interpreted wholly as methane leakage.

Public-power coal uses station consumption times rank-specific NCV. Generation and heat rates provide checks or explicit alternatives, not another input on top of the consumed fuel. K-Electric membership is tested fuel by fuel. Industrial coal distinguishes cement integrated kilns, brick firing, other industrial uses, captive fuel and unresolved supply/stock/boundary quantities. Dated fleet observations do not constitute an annual operating census.

IPPU uses source-specific physical output and trade with the correct process boundary. Cement/clinker production and trade, ammonia/feedstock and recovered carbon, steel and soda ash remain distinct from their combustion fuel. An aggregate manufacturing index does not supply missing physical clinker, ammonia or steel output. Missing process/control records remain coverage statuses.

\subsection{Household, agriculture and waste activity}
Household quantities combine deliveries and HIES survey evidence with explicit interpolation/extrapolation. The S dung reference is observed household use reconstructed between survey rounds; L assigns all ledger-diverted dung to household combustion. Dry solids including ash, volatile solids, as-received fuel and fuel energy are distinct quantities.

Crop area/production and fertiliser nutrients use compatible fiscal controls. Livestock is reconstructed consistently around the adopted census rather than spliced from an older projected baseline. The nitrogen ledger follows management, grazing and applied pools. Fertiliser offtake is distinct from application and release timing.

Waste advances the disposal-history state with demographic activity and the supported treatment/fate formulation. Fixed management parameters are not annual management observations. Waste origin, disposal, treatment and burning location are distinct supports. Aviation uses domestic traffic and fleet/LTO assumptions separately from fuel uplift, and railway traction/non-traction remains explicit.
